\documentclass[aspectratio=169]{beamer}

% Presentation settings
\usetheme{metropolis}
\usefonttheme[onlymath]{serif}

% Setting docuemnt language and fonts
\RequirePackage[bidi=basic, layout=tabular, provide=*]{babel}
\babelprovide[main, import]{hebrew}
\babelprovide{rl}
\babelfont[hebrew]{rm}[Renderer=Harfbuzz]{Libertinus Serif}
\babelfont[hebrew]{sf}{Liberation Sans}
\babelfont[hebrew]{tt}{Libertinus Mono}

\title{רקע מתמטי קצר על קבוצות ולוגיקה}
\author{}
\institute{80177 --- חשבון אינפיניטסימלי לתלמידי הנדסה ומדעים}
\date{}

\begin{document}

\frame{\titlepage}

\section{קבוצות}
\begin{frame}{המושג קבוצה}
	קבוצה היא אוסף של איברים. אנו מציינים קבוצה בעזרת סוגריים מסולסלים, שביניהם מגדירים מיהם האיברים של הקבוצה, על ידי כתיבה מפורשת שלהם בהפרדת פסיקים.

\end{frame}
\begin{frame}{דוגמה לקבוצה}
	$\{1,2,3\}$ היא הקבוצה שאיבריה הם המספרים $1$, $2$ ו־$3$.

	הקבוצה הזו לא מכילה את $5$, כי $5$ שונה מ־$1$, מ־$2$ ומ־$3$ והוא לא רשום בין הצומדיים (סוגריים מסולסלים).
\end{frame}


\begin{frame}{שייכות של איבר לקבוצה}
	הסימון $x\in A$ מציין שהאובייקט $x$ שייך לקבוצה $A$.

	לדוגמה $1 \in \{1,2,3\}$

	הסימון $y\notin A$ מציין שהאובייקט $y$ לא שייך לקבוצה $A$.

	לדוגמה $5\notin\{1,2,3\}$
\end{frame}

\begin{frame}{תיאור קבוצה באמצעות כלל}
	אפשר גם לתאר קבוצה בעזרת כלל. לדוגמה $\{ \text{המספרים השלמים בין $1$ ל־$3$} \} = \{1,2,3\}$

	הכלל שמגדיר מי שייך לקבוצה ומי לא חייב להיות חד־משמעי וברור. 

	יש סימון סטנדרטי כדי לתאר קבוצה שאיבריה נתונים באמצעות כלל המאופיין אותם.

	נדגים את הסימון על הקבוצה לעיל, אותה מסמנים $\{ x \mid 1 \le x \le 3, \text{$x$ מספר שלם} \}$
	נסביר בשלבים איך לקרוא את השורה הזו:
	\begin{itemize}
		\item ראשית זכרו – \textbf{קוראים את התיאור משמאל לימין!}
		\item שנית, הקו האנכי $\mid$ פירושו ''כך ש־''
	\end{itemize}
	לכן כתוב שם, ''הקבוצה של כל האובייקטים $x$ כך ש־$x$ מספר שלם ו־$x$ בין $1$ ל־$3$, כולל''.
\end{frame}

\begin{frame}{קבוצות מוכרות}
frameישנן קבוצות  אשר מופיעות כל הזמן במתמטיקה, ולכן נשמרו עבורן סימנים מוסכמים. נציג חלק מהם עכשיו
	\begin{enumerate}
		\item \textbf{מספר טבעי }הוא מספר שלם חיובי. קבוצת\textbf{ המספרים הטבעיים } היא הקבוצה $\mathbb{N}=\{1,2,3,4,\ldots\}$.
			למשל, $1\in\mathbb{N}$ ו־$10234\in\mathbb{N}$ אבל $0\notin\mathbb{N}$, $\:-1\notin\mathbb{N}$ וגם $\frac{3}{4}\notin\mathbb{N}$
		\item קבוצת ה\textbf{מספרים השלמים} היא הקבוצה $\mathbb{Z}=\{\ldots,-2,-1,0,-1,-2,\ldots\}$

			שימו לב: מספרים שלמים אינם בהכרח חיוביים! למשל, $\:-20\in\mathbb{Z}$
			וגם $2^{13}\in\mathbb{Z}$, אבל $\frac{3}{4}\notin\mathbb{Z}$
		\item \textbf{מספר רציונלי} הוא שבר, כלומר מספר מנה $\frac{m}{n}$ עבור
			$m$ ו־$n$ שלמים ו־$n\neq0$. מסמנים קבוצה זו על־ידי,
			\begin{center}
			$\begin{array}[t]{ccc}
			\mathbb{Q} & \underset{}{=} & \underset{}{\overset{}{\left\{ \left.\frac{m}{n}\right|n\neq0\text{ \:\text{ו־} }n\in\mathbb{Z}\text{ \:\text{ו־} }m\in\mathbb{Z}\right\} }}\end{array}$  
			\par\end{center}
	\end{enumerate}
\end{frame}

\begin{frame}{שוויון בין קבוצות}
	שתי קבוצות הן \textbf{שוות} אם יש להן את אותם האיברים.

	כלומר, כל איבר של הקבוצה הראשונה שייך לשניה, וכל איבר של הקבוצה השניה שייך לראשונה.
\end{frame}



\begin{frame}{הכלה ותת־קבוצות}
	אומרים שקבוצה $A$ \textbf{מוכלת }בקבוצה $B$ אם כל איבר של
	$A$ הוא גם איבר של $B$. 
	במקרה זה אומרים גם ש־$A$ היא \textbf{תת־קבוצה} של $B$.

	מסמנים זאת $A\subseteq B$.

	לדוגמה, $\{5,11,13\}\subseteq\{2,5,11,12,13,17\}$.\bigskip{}

	\textbf{הערה}: שימו לב שהסמל $\subseteq$ מקשר בין שתי קבוצות והסמל $\in$ מקשר בין קבוצה ואיבר שלה.
\end{frame}

\begin{frame}{דוגמה}
	תהיינה $A=\{1,2\},B=\{2,3\}$.

	אז $A\not\subseteq B$ כי $1\in A$ אבל $1 \notin B$.

	באופן דומה $B\not\subseteq A$ כי $3\in B$ אבל $3\notin A$.
\end{frame}


\begin{frame}{שוויון והכלה}
	קודם אמרנו שקבוצות שוות אם יש להן אותם איברים, לכן בהתאם נקבל $A = B$ אם ורק אם $A \subseteq B$ וגם $B \subseteq A$.

	לעיקרון זה נקרא הכלה דו־כיוונית.
\end{frame}

\section{השפה המתמטית}

\begin{frame}{פסוקים}
	היחידה הבסיסית של הלוגיקה היא \textbf{פסוק}. פסוק הוא אמירה שניתן
	לשייך לה ערך '' אמת'' או '' שקר''{}
	בצורה חד משמעית. 

	\textbf{דוגמה}: $0=1$ זהו פסוק שקר.

	לא כל מה שכותבים, גם בסימנים מתמטיים, הוא פסוק. למשל ''$x^{4}\ge0$'' איננו פסוק.

	הטענה לא כתובה באופן תקין --- לא נאמר בה מיהו $x$. אי אפשר ליחס לאמירה הזו ערך אמת בלי שיודעים מיהו $x$. יתר על כן, יש גם בחירות של $x$ עבורם לביטוי ''$x^{4}\ge0$'' אין שום משמעות,
	למשל אם $x$ היא מכונית. לכן כדי להפוך את זה לפסוק יש להבהיר מיהו $x$.

	\textbf{טענה}: לכל $x\in\mathbb{Q}$ מתקיים $x^{4}\ge0$

	זו טענה שמנוסחת באופן תקין!
\end{frame}

\begin{frame}{כמתים}
frameמהאמור לעיל נובע ש־$x>0$ איננו פסוק. אפשר להפוך את הביטוי $x>0$
	לפסוק על ידי \textbf{כימות} על $x$, 
	\begin{itemize}
		\item עם \textbf{הכמת האוניברסלי $\left(\forall \right)$}: למשל $x>0$ עבור כל $x$ שלם (בסמלים  $\forall x\in\mathbb{Z}\quad x>0$). 
			במילים אחרות זהו הפסוק ”כל מספר שלם הינו חיובי''.
		\item עם\textbf{ הכמת הקיומי $\left(\exists \right)$}: למשל  קיים $x$
			רציונלי כך ש־$x>0$ (בסמלים  $\exists x\in\mathbb{Q}\quad x>0$)
	\end{itemize}
\end{frame}

\begin{frame}
	באופן כללי תהי $P(x)$ אמירה על משתנה $x$. אפשר להפוך את האמירה
	לפסוק שאיננו תלוי ב־$x$ על ידי כמתים:

	בהינתן קבוצה $S$, אפשר ליצור את הפסוקים,
	\begin{itemize}
		\item לכל $x\in S$ מתקיים $P(x)$. בסימנים, $\forall x\in S\ P(x)$
		\item קיים $x\in S$ עבורו $P(x)$ מתקיים. בסימנים: $\exists x\in S\ P(x)$
	\end{itemize}
\end{frame}

\begin{frame}{שלילת כמתים}
	\textbf{דוגמה, }התבוננו בפסוק הבא: לכל $x\in\mathbb{Q}$ מתקיים
	$x\cdot x=x+x$.\smallskip{}

	\fcolorbox{black}{white}{\begin{minipage}{\textwidth}
		כדי להראות שלא נכון ש\textbf{כל איברי }הקבוצה $A$ מקיימים תכונה
		מסויימת, צריך להראות שיש לפחות \textbf{איבר אחד} ב־$A$ שלא מקיים
		זאת. כלומר, יש לתת \textbf{דוגמה נגדית}.
	\end{minipage}}

	שלילת הפסוק בדוגמה שקולה לטענה הבאה\textbf{ }: קיים $x\in\mathbb{Q}$
	כך ש־$x\cdot x\neq x+x$.

	\textbf{הוכחה}: $x=1\in\mathbb{Q}$ מקיים $1\cdot1=1\neq2=1+1$. 

	באופן כללי, תהי $P(x)$ אמירה על משתנה $x$ ו־$A$ קבוצת
	איברים. השלילה של הפסוק $\forall x\in AP(x)$ היא $\exists x\in A\neg P(x)$. 

	כלומר: $\neg(\:\forall x\in AP(x)\:)\quad\iff\quad\exists x\in A\neg P(x)$.
\end{frame}

\begin{frame}{שלילת כמתים}
	\textbf{דוגמה :} התבוננו בפסוק הבא: קיים $x\in\mathbb{Q}$
	כך ש־$x+1\leqslant x$.\smallskip{}

	\fcolorbox{black}{white}{\begin{minipage}{\textwidth}
		כדי להראות ש\textbf{לא קיים} $x$ עם תכונה מסוימת, צריך להראות ש\textbf{כל $x$ }מקיים את שלילת התכונה.
	\end{minipage}}

	כלומר, שלילת הטענה הקודמת שקולה לטענה הבאה: כל $x\in\mathbb{Q}$
	מתקיים $x<x+1$.

	\textbf{הוכחה}: יהי $x\in\mathbb{Q}$. נוסיף את $x$ לשני האגפים של אי־השוויון $0<1$ ונקבל $x<x+1$.

	באופן כללי, $\neg(\:\exists x\in AP(x)\:)\quad\iff\quad\forall x\in A\neg P(x)$
\end{frame}

\begin{frame}{סדר כמתים}
	\fcolorbox{black}{white}{\begin{minipage}{0.9 \textwidth}
		כשיש יותר מכמת אחד, \textbf{הסדר בו הם כתובים הוא קריטי}, ושינוי הסדר יכול לשנות לחלוטין את משמעות הפסוק.
	\end{minipage}}

	התבוננו למשל בפסוקים הבאים,
	\begin{itemize}
		\item $\forall x\in\mathbb{Q}\quad\exists y\in\mathbb{Q}\ y>x$
		\item $\exists y\in\mathbb{Q}\quad\forall x\in\mathbb{Q}\ y>x$
	\end{itemize}


	הפסוק הראשון אומר שעבור כל רציונלי $x$ יש רציונלי $y$ הגדול ממנו. ראינו קודם שהטענה הזו נכונה.

	לעומת זאת הפסוק השני אומר שקיים רציונלי $y$ כך שכל איברי $\mathbb{Q}$ קטנים ממנו. זה בוודאי לא נכון כי למשל $y$ איננו גדול מעצמו, ולכן הטענה נכשלת עבור $x=y$ (וגם עבור $x=y+1$ למשל).
\end{frame}

\end{document}
